% -*- TeX-command-extra-options: "-shell-escape" -*-
% This file requires --shell-escape (for \write18 calls).
% The magic comment above makes AUCTeX (Emacs) pass --shell-escape automatically.
% Command-line users: pdflatex --shell-escape CacheTest.tex
%
% CacheTest.tex — test suite for smart checksum-based cache invalidation (Issue #10)
%                 and \inln fallback when output file is missing (Issue #16)
%
% HOW TO VERIFY CACHING:
%   1st compile: all code runs; generated/ gets .md5 sidecar files.
%   2nd compile: no code reruns; .tex output file mtimes are unchanged.
%   Edit one chunk: only that chunk reruns on next compile.
%   [run] option: forces all code to rerun regardless of checksums.
%   [cache] option: never reruns; Issue #16 fix — runs anyway if output missing.
%
% Languages exercised: R, Python, Julia

\documentclass[12pt]{article}
\usepackage[nohup,R,python,julia,stopserver,listings]{../../runcode}
\usepackage[margin=1in]{geometry}

\begin{document}

\title{\texttt{CacheTest}: Smart Cache and \texttt{{\textbackslash}inln} Fallback}
\date{}
\maketitle

% ─────────────────────────────────────────────────────────────────────────────
\section{R}

\subsection{\texttt{\textbackslash runRChunk} — per-chunk cache via MD5}

The chunks below are cached individually: editing \texttt{r-sum} will only
rerun that chunk on the next compile, leaving \texttt{r-mean} untouched.

Sum:  \runRChunk{code/hello.R}{r-sum}
Mean: \runRChunk{code/hello.R}{r-mean}

\subsection{\texttt{\textbackslash runRIncOut} — whole-file cache}

Running the full file produces all three outputs:

\runRIncOut{code/hello.R}

\subsection{\texttt{\textbackslash inlnR} — backtick inline (Issue \#16 fallback)}

On first compile there is no cached \texttt{.tex} output, so the code must run
regardless of the \texttt{\textbackslash ifruncode} flag.

Inline sum: \inlnR{```cat(sum(c(2,4,6,8,10)))```}

% ─────────────────────────────────────────────────────────────────────────────
\section{Python}

\subsection{\texttt{\textbackslash runPythonChunk} — per-chunk cache via MD5}

Sum:  \runPythonChunk{code/hello.py}{py-sum}
Mean: \runPythonChunk{code/hello.py}{py-mean}

\subsection{\texttt{\textbackslash runPythonIncOut} — whole-file cache}

\runPythonIncOut{code/hello.py}

\subsection{\texttt{\textbackslash inlnPython} — backtick inline (Issue \#16 fallback)}

Inline sum: \inlnPython{```print(sum([2,4,6,8,10]))```}

% ─────────────────────────────────────────────────────────────────────────────
\section{Julia}

\subsection{\texttt{\textbackslash runJuliaChunk} — per-chunk cache via MD5}

Sum:  \runJuliaChunk{code/hello.jl}{jl-sum}
Mean: \runJuliaChunk{code/hello.jl}{jl-mean}

\subsection{\texttt{\textbackslash runJuliaIncOut} — whole-file cache}

\runJuliaIncOut{code/hello.jl}

\subsection{\texttt{\textbackslash inlnJulia} — backtick inline (Issue \#16 fallback)}

Inline sum: \inlnJulia{```println(sum([2,4,6,8,10]))```}

% ─────────────────────────────────────────────────────────────────────────────
\section{Issue \#16: \texttt{[cache]} package option with missing output}

Compiled with \texttt{[cache]} the package sets \texttt{\textbackslash runcodefalse}.
The Issue \#16 fix ensures that if the output \texttt{.tex} file does not exist,
the code still runs.  Delete \texttt{generated/} and recompile with
\texttt{\textbackslash usepackage[cache,...]\{runcode\}} to exercise this path.

Fallback: \inlnR{```cat("cache-fallback")```}[inlnR-fallback]

\end{document}
